\documentclass[12pt]{amsart}
\pagestyle{empty}
\usepackage[hmargin=1in,vmargin=1in]{geometry}
\usepackage{url,  mathrsfs}
\usepackage{hyperref}
\usepackage{fancyhdr} 
\usepackage[T1]{fontenc}

\usepackage{listings,xcolor}

\lstset{language=Mathematica}
\lstset{basicstyle={\sffamily\footnotesize},
  numbers=left,
  numberstyle=\tiny\color{gray},
  numbersep=5pt,
  breaklines=true,
  captionpos={t},
  frame={lines},
  rulecolor=\color{black},
  framerule=0.5pt,
  columns=flexible,
  tabsize=2
}

\usepackage{url,  mathrsfs}
\usepackage{hyperref}
\usepackage{amsmath}
\usepackage{cancel}
\usepackage{amssymb}
\usepackage{graphicx,  epstopdf}
\usepackage{subfig}
\usepackage{color}
\usepackage{bbm}
\usepackage{subfloat}
\usepackage[draft]{fixme}
\usepackage{bm} 
\usepackage{bbm} 
\usepackage{epigraph}
\usepackage{endnotes}
\usepackage{ifpdf}
\usepackage{array}

\usepackage{multicol}
\usepackage{multirow}
\usepackage{graphicx}
\usepackage{tikz}

\newcommand{\A}{\mathbb{A}}
\newcommand{\Q}{\mathbb{Q}}
\newcommand{\N}{\mathbb{N}}
\newcommand{\Z}{\mathbb{Z}}
\newcommand{\R}{\mathbb{R}}
\newcommand{\F}{\mathbb{F}}
\newcommand{\C}{\mathbb{C}}
\renewcommand{\P}{\mathbb{P}}
\newcommand{\I}{\mathcal{I}}
\newcommand{\V}{\mathcal{V}}
\newcommand{\cH}{\mathcal{H}}
\newcommand{\ol}{\overline}
\newcommand{\J}{\mathcal{J}}
\newcommand{\X}{\mathcal{X}}
\theoremstyle{definition}
\newtheorem{prob}{Problem}
\newtheorem{extraProb}{Extra Problem}

\DeclareMathOperator{\conv}{conv}
\DeclareMathOperator{\cone}{cone}
\DeclareMathOperator{\ext}{ext}




\begin{document}
 \thispagestyle{empty}
\begin{center} {\large \textbf{Math 380 -- Homework  5}} \smallskip \\ 
Due on Thursday, May 14\ \\ \ \\
 \end{center}

You are welcome to talk with other students in the class about problems but should write up solutions on your own. 
Solutions can be handwritten or typed but need to be legible and submitted via Gradescope by the end of the day 
on Thursday. You should justify all your answers in order to receive full credit.  
\bigskip 

\begin{flushleft}

{\bf Good practice problems} (not to be turned in): \\

CLO Section 3.1, Exercises 4, 5, 6 \\
CLO Section 3.2, Exercises 3, 4 \\
CLO Section 3.3, Exercises 2, 3, 6, 10, 11, 13, 14\\
\ \ \\
\bigskip 




\begin{prob} CLO \S3.1, Exercise 7 \smallskip \\
\end{prob}
\bigskip 

Consider the ideal 
\[ I = \langle t^2 + x^2 + y^2 + z^2 , t^2 + 2x^2 - xy - z^2, t+y^3 - z^3\rangle \subseteq \mathbb Q[t,x,y,z]. \]

\begin{enumerate}
    \item[(a)] With the lexicographical order $t > x > y > z$, we find a Gr\"obner basis for $I$:
    \begin{lstlisting}[language=Mathematica]
In[1]:= G1 = 
 GroebnerBasis[{t^2 + x^2 + y^2 + z^2, t^2 + 2 x^2 - x y - z^2, 
    t + y^3 - z^3}, {t, x, y, z}] // FullSimplify

Out[1]= {5 y^4 + 5 y^8 + y^12 - 10 y^5 z^3 - 4 y^9 z^3 - 
  4 y^3 z^5 (3 + z^4) + z^4 (3 + z^4)^2 + y^2 z^2 (13 + 5 z^4) + 
  6 y^6 (z^2 + z^6), -y^3 (5 + 5 y^4 + y^8) + (3 x - 7 y - 
     3 y^5) z^2 + 2 y^4 (5 + 2 y^4) z^3 + 
  6 y^2 z^5 + (x - 3 y - 5 y^5) z^6 + 2 y^2 z^9, 
 y (x + 2 y + y^5) + 3 z^2 - 2 y^3 z^3 + z^6, 
 x^2 + y^2 + y^6 + z^2 - 2 y^3 z^3 + z^6, t + y^3 - z^3}

In[2]:= (* calculate total degrees *)
  Max[Total /@ Keys[CoefficientRules[#, {t, x, y, z}]]] & /@ G1 

Out[2]= {12, 11, 6, 6, 3}
    \end{lstlisting}
    Eliminating $t$, the first elimination ideal $I_1 = \langle G_1\rangle= I \cap \mathbb Q[x,y,z]$ is generated by the first four polynomials in the output list \verb|Out[1]|.

    \item[(b)] We can compute a Gr\"obner basis for $I_1$ using grevlex order:
    \begin{lstlisting}[language=Mathematica]
In[3]:= GroebnerBasis[G1[[;; -2]], {x, y, z}, 
 MonomialOrder -> DegreeReverseLexicographic]

Out[3]= {x^2 - x y - y^2 - 2 z^2, 
 x y + 2 y^2 + y^6 + 3 z^2 - 2 y^3 z^3 + z^6}
    \end{lstlisting}
    \item[(c)] We can fix the $>_{l=1}$ elimination order, where $\alpha > \beta$ if and only if $\alpha_1 > \beta_1$ or $\alpha_1 = \beta_1$ and $\alpha >_\text{grevlex} \beta$. Take the following generators:
    \begin{align*}
    f_1 = x^2-xy-y^2-2z^2 &\quad f_2 =xy+2y^2+y^6+3z^2-2y^3z^3+z^6 \\ f_3=t+y^3-z^3.
    \end{align*}
    Note that in part (b) we found that $I_1 =\langle f_1, f_2\rangle$ where we consider $I_1$ as an ideal of $\mathbb Q[x,y,z]$. If we consider $\langle I_1\rangle$ as the ideal of $\mathbb Q[t,x,y,z]$ generated by $I_1$, part (a) shows that $I = \langle I_1 \rangle + \langle t+y^3-z^3\rangle$ since every other generator listed is a member of $I_1$. Thus, since $\langle I_1\rangle = \langle f_1,f_2\rangle$, we have $I = \langle f_1,f_2,f_3\rangle$.
    
    Let $f \in I$ be any polynomial. If $f \in I_1$, then $f \in \mathbb Q[x,y,z]$, and the $>_{l=1}$ elimination order restricted to this subset behaves exactly like the $>_\text{grevlex}$ order. Since $G_0=\{ f_1,f_2\}$ is a Gr\"obner basis for $I_1$ under the $>_\text{grevlex}$ order, we know that \[\text{LT}_{l=1}(f) =\text{LT}_\text{grevlex}(f) \in \langle \text{LT}_\text{grevlex}(G_0)\rangle = \langle \text{LT}_{l=1}(G_0)\rangle.\] Here, we are considering ideals of $\mathbb Q[x,y,z]$. On the other hand, if $f \not \in I_1$, then some term of $f$ must involve $t$, and then $\text{LT}_{l=1}(f)$ is simply the term with the highest degree in $t$. Then, $\text{LT}_{l=1}(f_3) = t$ must divide $\text{LT}_{l=1}(f)=t^m \cdot h$ for $m \ge 1$ and $h \in \mathbb Q[x,y,z]$. Thus, $G = \{f_1,f_2,f_3\}$ is a Gr\"obner basis for $I$ under $>_{l=1}$ since we have shown that $\text{LT}_{l=1}(f_i)$ for some $f_i \in G$ must divide $\text{LT}_{l=1}(f)$ for all $f \in I$.

    \begin{lstlisting}[language=Mathematica]
In[4]:= G1ast = 
 GroebnerBasis[{t^2 + x^2 + y^2 + z^2, t^2 + 2 x^2 - x y - z^2, 
   t + y^3 - z^3}, {t, x, y, z}, MonomialOrder -> EliminationOrder]

Out[4]= {x^2 - x y - y^2 - 2 z^2, t^2 + x y + 2 y^2 + 3 z^2, 
 t + y^3 - z^3}

In[5]:= G1ast[[2]] /. {t -> z^3 - y^3} // ExpandAll

Out[5]= x y + 2 y^2 + y^6 + 3 z^2 - 2 y^3 z^3 + z^6
    \end{lstlisting}
\end{enumerate}

\begin{prob} CLO \S3.2, Exercise 5\\
\end{prob}
\bigskip 
Let $I \subseteq \mathbb C[x,y]$ be an ideal where $I_1 \neq \langle 0 \rangle$. Since $\mathbb C[y]$ is a principal ideal domain, we can write $I_1 = \langle f\rangle \subseteq \mathbb C[y]$. If $f$ is a constant polynomial, then $\mathbf V(I_1)=\varnothing$ since $f \neq 0$, and since $\pi_1(\mathbf V(I))\subseteq \mathbf V(I_1)$, we get $\pi_1(\mathbf V(I))= \mathbf V(I_1)=\varnothing$. Otherwise, since $f$ is a non-constant polynomial, it must have at most $\deg f$ roots, so $\mathbf V(I_1)$ is a finite set with cardinality at most $\deg f$. Then, $\pi_1(\mathbf V(I))\subseteq \mathbf V(I_1)$ is also a finite set. Since all finite sets are varieties, there exists a variety exactly equal to $\pi_1 (\mathbf V(I))$, and we get $\pi_1(\mathbf V(I))=\mathbf V(I_1)$ as $\mathbf V(I_1)$ is the smallest variety containing $\pi_1(\mathbf V(I))$ by the Closure theorem ($\mathbb C$ is algebraically closed).

\newpage{}
\begin{prob} CLO \S3.3, Exercise 7 \smallskip
 \\

Consider the following parametric surface $S \subseteq\mathbb C^3$:
\begin{align*}
    x &= uv, \\
    y &= uv^2, \\
    z &= u^2.
\end{align*}
\begin{enumerate}
    \item[(a)] Since $\mathbb C$ is algebraically closed, we use the Closure theorem to find the smallest variety containing $S$, finding a Gr\"obner basis for $I=\langle x-uv,y-uv^2,z-u^2\rangle$ under the $u > v>x>y>z$ lexicographic order and computing the second elimination ideal.
    \begin{lstlisting}[language=Mathematica]
In[6]:= G3 = 
 GroebnerBasis[{x - u v, y - u v^2, z - u^2}, {u, v, x, y, z}, 
  MonomialOrder -> Lexicographic]

Out[6]= {x^4 - y^2 z, -x^3 + v y z, 
 v x - y, -x^2 + v^2 z, -x^2 + u y, u x - v z, u v - x, u^2 - z}
    \end{lstlisting}
    Since $I_2 = \langle G_2\rangle$, we get $I_2 = \langle x^4-y^2 z\rangle$. Thus, $V=\mathbf V(x^4-y^2z)$ is the smallest variety in $\mathbb C^3$ that contains $S$.
    \item[(b)] First, we can show that $\mathbf V(I_1) = \pi_1(\mathbf V(I))$. We can write the generators of $I$ in the following form:
    \[ 
        I = \langle 0\cdot u +x^4-y^2z, \dots, x \cdot u-vz, v\cdot u - x, 1 \cdot u^2 - z \rangle.
    \]
    By the geometric extension theorem, we have that
    \[ \mathbf V(I_1) = \pi_1(\mathbf V(I))\cup [\mathbf V(I_1)\cap \mathbf V(0,\dots, x,v,1)].\]
    Since $\mathbf V(0,\dots, x,v,1) = \varnothing$, the right hand side is the empty set and $\mathbf V(I_1) = \pi_1(\mathbf V(I))$. Note that this means $S = \pi_2 (\mathbf V(I)) = \pi_1 (\pi_1 (\mathbf V(I)) = \pi_1 (\mathbf V(I_1))$.
    
    Now, we can apply the extension theorem again. We have \[I_1 = \langle x^4 -y^2 z, v y z - x^3, vx-y, v^2 z - x^2 \rangle,\] and $I_2$ is the first elimination ideal of $I_1$. We can rewrite the generators of $I_1$:
    \[
    I_1 = \langle 0 \cdot v+x^4 -y^2, \; yz\cdot v-x^3,\; x\cdot v-y, \; z\cdot v^2-x^2\rangle.
    \]
    Thus, by the geometric extension theorem, \[V=\mathbf V(I_2)=\pi_1(\mathbf V(I_1)) \cup [\mathbf V(I_2)\cap \mathbf V(yz,x,z)].\]
    Again note that $S = \pi_1(\mathbf V(I_1))$. This hints that points in $V$ that are not in $S$ are of the form $(0,y,0)$, and we can verify that all points of that form satisfy the equation for $V$. However, $z=0$ implies that $u=0$ and so $x=y=z=0$, so $(0,0,0)$ is the only such point in $S$. We have found all the points in $V$ that do not lie in $S$, and $V \setminus S = \{(0,y,0) \in \mathbb C^3 : y \neq 0\}$.
\end{enumerate}

\end{prob}
\bigskip 




\end{flushleft}
\end{document}